
\begin{proposition}[Ordinary-triple images have arithmetic genus at most one]\label{cdp:st:A.5.19}
Let $f:Y\to C_0$ be a proper holomorphic map from a smooth
compact complex threefold with $b_2(Y)=0$, with connected
fibres and smooth general fibre a complex two-torus.
Let $mS$ be a full fibre with irreducible reduced support,
and let $D\subset\operatorname{Sing}S$ be an irreducible
curve with generic ordinary-triple plane transversal.
Then $D$ is a local complete intersection in $Y$ and
$p_a(D)\le1$. The possibilities for its underlying curve
are a smooth rational curve, a rational curve with one
ordinary node or cusp, and a smooth elliptic curve.
In the elliptic case the transversal type is constant.

When $D$ is singular rational, every component of the
normalized conductor curve dominating $D$ is rational,
and its map to $D^\nu\simeq\PP^1$ has degree one.
Consequently, in the irrational ruled setup of \ref{cdp:st:A.5.18}
with the generic transversal replaced by an ordinary
triple point, the image of the noncontracted horizontal
section is either smooth rational or smooth elliptic.
In the elliptic case the horizontal cover is unramified
of degree three, the normalization is smooth minimal
elliptic ruled, and the fibre has Euler number zero.
In the rational case the horizontal cover has degree
two or three; these cases are not excluded by this
statement.
\end{proposition}
\begin{proof}
By \ref{cdp:st:A.5.14}, $p_a(D)\le2$. If its normalization has
positive genus, that statement and \ref{cdp:st:A.5.5} already give
a smooth elliptic curve of constant transverse type.
Suppose its normalization $\nu:\Sigma\to D$ is
$\PP^1$, with total delta invariant $\delta\le2$.

First let $D$ be lci. Put
$N=(\nu^*(I_D/I_D^2))^*$ and
$M=\nu^*(\mathcal O_Y(S)|_D)$.
The quadratic normal jet vanishes generically, so
its vanishing descends on the reduced curve through
the locally free graded piece $I_D^2/I_D^3$.
Thus the initial cubic is an actual section of
$\operatorname{Sym}^3N^*\otimes M$.
Adjunction and the curve conductor give
$\deg N=-2+2\delta$, while $\deg M=0$.
Its nonzero binary-cubic discriminant is a section
of $M^4\otimes(\det N)^{-6}$, of degree
\[
                         12-12\delta.             \tag{A.U.1}\label{cdp:eq:A.U.1}
\]
Hence $\delta\le1$.

We exclude every non-lci germ of delta at most two
by its symbolic cube. The one- and two-branch analytic
classification is proved in (A.T.14). There is one
additional possibility when three branches occur:
delta at least two is forced by identifying their
constants, so equality makes the ring exactly
\[
 \{(a,b,c)\in\C\{t_1\}\oplus\C\{t_2\}\oplus\C\{t_3\}:
                      a(0)=b(0)=c(0)\}.
\]
Its actual embedding in the smooth threefold is the
union of the three coordinate axes, with ideal
$I=(xy,xz,yz)$. Indeed the three branch parameters
give a basis of the cotangent space, so their lifts
are ambient coordinates as in \ref{cdp:st:A.5.18}.
Thus there are exactly three non-lci cases to consider;
each consumes the whole delta budget at one point.

For the space cusp retain $I=(g_1,g_2,g_3)$ and
$F$ from (A.T.15). Its exact symbolic-cube identity is
\[
                         I^{(3)}=I^3+FI.           \tag{A.U.2}\label{cdp:eq:A.U.2}
\]
Here is a verification in the convergent local ring.
The right-hand side agrees with $I^3$ after inverting
$x$, by (A.T.15). Its quotient is finite over
$\C\{x\}$, of generic rank $3\cdot6=18$: a square
normal plane has six monomials of total degree less
than three. Modulo $x$, its ideal in $\C\{y,z\}$ is
\[
 (y,z)^6+(y^2z^3,yz^4,z^5).
\]
The fifteen monomials of total degree at most four,
together with $y^5,y^4z,y^3z^2$, give residue dimension
eighteen. The same discrete-valuation-module argument
used for (A.T.15) makes the quotient $x$-torsion-free.
It is therefore the saturation of $I^3$, which equals
the symbolic cube since the curve is lci off the origin.
This proves \eqref{cdp:eq:A.U.2}.

For the cusp-plus-line ideal $I=(xz,yz,G)$ with
$G=y^2-x^3$, one has
\[
                   I^{(3)}=I^3+(zG)I.             \tag{A.U.3}\label{cdp:eq:A.U.3}
\]
To verify this directly, write $P=(z,G)$ and
$Q=(x,y)$. Their powers are primary because these
prime ideals are complete intersections. The
symbolic cube is $P^3\cap Q^3$.
Write an element of $P^3$ as
$z^3A+z^2GB+zG^2C+G^3E$.
The last two terms are already in $Q^3$.
Reducing the other terms modulo $Q^3$ gives
$zA+GB=0$ there. The constant-in-$(x,y)$ coefficient
of $B$ must therefore be divisible by $z$; write
$B=zB'+B_Q$, $B_Q\in Q$.
It follows that $A+GB'\in Q^3$.
Consequently
\[
 P^3\cap Q^3=(z^3Q^3,z^2GQ,zG^2,G^3)
                        =I^3+(zG)I,
\]
as claimed.

For the three axes, membership of
$x^a y^b z^c$ in the symbolic cube means
$a+b,a+c,b+c\ge3$.
If one exponent is zero, the other two are at
least three and the monomial belongs to $I^3$.
If all are positive, either at least two are at
least two, in which case it belongs to $(xyz)I$,
or two are one, which violates one of the three
inequalities. Hence
\[
                         I^{(3)}=I^3+(xyz)I.       \tag{A.U.4}\label{cdp:eq:A.U.4}
\]

In all three cases the generic ordinary-triple
equation $f_0$ belongs to $I^{(3)}$: on the smooth
locus of each curve branch its order in the two
normal variables is at least three. Define the
bundle by its actual conormal image,
\[
 K=\operatorname{im}\bigl(\nu^*(I/I^2)
                           \to\nu^*\Omega_Y^1\bigr),
 \qquad N=K^*.
\]
For the first two cases (A.T.18) gives $\deg N=4$.
For the three axes, on the branch $(t,0,0)$ the
frame is $t\,dy,t\,dz$, with conormal determinant
order two and no tangent ramification. Each of the
three branches has the same correction, for total
six, so again $\deg N=-2+6=4$.

Use the actual affine maps (A.T.20)--(A.T.21).
For the three axes add
\[
              (u,v,t)\longmapsto(t,u/t,v/t)
                                                        \tag{A.U.5}\label{cdp:eq:A.U.5}
\]
on the first branch, and the two analogous permuted
maps on the other branches. In these charts the
first normal derivatives of the ideal generators
are holomorphic in the specified conormal frame.
An $I^3$ term consequently has holomorphic initial
cubic. The extra term in each of \eqref{cdp:eq:A.U.2}--\eqref{cdp:eq:A.U.4}
has initial cubic equal to the quadratic term of
its displayed symbolic-square generator, multiplied
by the first derivative of an ideal generator.

For the space cusp that quadratic term is (A.T.16),
so the cubic pole is at most three. For the
cusp-plus-line germ it is $t^{-2}uv$ on the cusp
and $s^{-1}v^2$ on the smooth branch, so the pole
bounds are two and one. For the three axes it is
$t^{-1}uv$ in \eqref{cdp:eq:A.U.5}, giving pole bound one on
each branch. Coefficients multiplying these terms
are holomorphic; their derivatives do not enter
the initial cubic because the lower normal jets
of the products vanish.

The initial cubic is intrinsic on the dense smooth
part of $D$, as all lower jets are zero. Its
transition under a normal-coordinate change is
therefore the actual cubic tensor in
$\operatorname{Sym}^3N^*\otimes M$; higher-order
coordinate terms make no contribution. The binary
discriminant is homogeneous of degree four in
these coefficients. It is a meromorphic section
of $M^4\otimes(\det N)^{-6}$, of degree $-24$,
with pole divisor of total degree at most twelve
in each of the three cases. After this exact pole
correction its line has negative degree, at most
$-12$. It cannot have the required nonzero
holomorphic section. This excludes all non-lci
cases and proves that $D$ is lci with $p_a(D)\le1$.
The list of underlying curves follows from the
delta-one classification in \ref{cdp:st:A.5.14}.

We give the finite-cover comparison at a singular
rational image. Here $\delta=1$, so \eqref{cdp:eq:A.U.1} has
degree zero and its nonzero section is nowhere zero
on $\Sigma=\PP^1$. The cubic defines a subscheme
$R_3\subset\PP(N)$ of tangent-root lines. Every
fibre consists of three distinct points. In a
projective coordinate chart, its derivative with
respect to the root coordinate is nonzero, so the
implicit function theorem makes
$R_3\to\Sigma$ a finite etale map of degree three.
There is no component in a whole projective fibre,
since its nonzero discriminant forbids an identically
zero cubic there.

Over the generic point of $D$, normalization of an
ordinary triple transversal separates exactly its
three tangent lines. This supplies an isomorphism
between the function-field algebra of $R_3$ and
that of the normalized conductor curves dominating
$D$, after their maps are factored through $\Sigma$.
Both are normal finite curves over $\Sigma$.
Taking the integral closure of $\mathcal O_\Sigma$
in this same algebra therefore identifies the two
finite covers everywhere. Every connected finite
etale cover of the simply connected Riemann sphere
has degree one. Thus each of these conductor
components is rational and maps isomorphically to
$\Sigma$, proving the second assertion.

In the irrational ruled setting the horizontal
section has positive genus. Its finite map to a
rational image normalization must have degree at
least two. The preceding comparison excludes a
singular rational image, and the generic triple
transversal bounds its degree above by three.
The remaining rational image is therefore smooth,
with degree two or three. If the image is elliptic,
no vertical rational curve can dominate it, so all
three branches belong to the horizontal section.
Statement \ref{cdp:st:C2d} gives an unramified degree-three
cover, normal degree zero and constant type.
The exact minimal-resolution factorization and
blowdown argument (A.T.22) applies without change
to this coefficient-two section: its proof uses
only those properties. It makes the normalization
smooth minimal elliptic ruled with conductor twice
the section. The actual normalization Euler formula
(A.T.23) then gives zero.
\end{proof}

The singular rational possibilities cannot be
removed by a blanket local spherical-cycle argument.
For example
\[
                     f=x(yz)(x-yz)
                                                        \tag{A.U.6}\label{cdp:eq:A.U.6}
\]
has ordinary-triple transversals along the nodal curve
$D=(x,yz)=0$ and an additional ordinary-double
critical line $y=z=0$. For $t\ne0$, put
$B_t=\{(x,v):xv(x-v)=t\}$. The exact maps
\[
 B_t\times\C^*\longrightarrow f^{-1}(t),
 \quad (x,v,y)\longmapsto(x,y,v/y),\qquad
 (x,y,z)\longmapsto(x,yz,y)
                                                        \tag{A.U.7}\label{cdp:eq:A.U.7}
\]
are inverse holomorphic maps. Nonzero $t$ forces
$v,y$ and $z$ to be nonzero. The smooth affine
curve $B_t$ and $\C^*$ are aspherical, so this
regular affine fibre has no spherical homotopy
class. This is a local model, not a compact
$b_2=0$ example. A deformation smoothing $yz$
changes the additional critical line's transverse
type on the boundary and does not satisfy the
comparison used in \ref{cdp:st:A.5.17}.

Exact ideal, normal-coordinate and discriminant
calculations are retained in
\path{code/E_triple_images.py} and
\path{code/E_triple_images.out}; the independently
computed symbolic saturations are also in
\path{code/A_low_delta_cusp.py}.


